% Transcription — fonds Alexandre Grothendieck, shelfmark 46, batch 4 (pages 61-80).
% Established from the facsimile archives/batches/46/batch-04.pdf (cut from a
% copy of 46.pdf fetched through the project relay,
% https://grothendieck-relay.onrender.com/source/46.pdf; PDF page k =
% archivists' page 60+k, no cover sheet), per the skill transcribe-grothendieck.
% The folder is ninety-four pages in five batches; this is the fourth.
%
% Pass: Opus 5.5 (claude-opus-5-5), 2026-09-26 — first pass, unchecked against the pages by a human.
%
% Pages skipped, and why: none as unrelated. Every leaf of 61-80 is in his
% hand. Two are covers (see below) and get no \page marker.
% Covers, no \page marker, words as headings with a \note: 67 (tan folder,
% « "Restriction des scalaires" dans les groupes algébriques et leurs espaces
% principaux homogènes... »), 71 (grey folder, « Petits fourbis
% cohomologiques », the word « fourbis » flagged as doubtful in the note).
%
% Typescript: none. No further leaf of the « Définition du morphisme
% Verschiebung » typescript of batch 2 (pages 23-26) occurs in pages 61-80,
% and no other typed leaf either; nothing here is addressed to him. So there
% is no typescript decision to make in this batch.
%
% Pages summarised: none. Page 66 carries three lines only. Two blocks
% cancelled by a large scribble (page 65 foot, page 75 middle) are given in a
% \note rather than re-set. The two bibliographical lines heading page 78
% (Peterson; Moore), struck through by a long wavy stroke, are recorded in a
% \note.
%
% His pagination: none. His numbered items: Questions 1-2 (page 64); on pages
% 76-79 « Lemme 1 » (page 76, finite-field lemma) and a new run « Lemme 1 » to
% « Lemme 5 » (pages 77-79, cohomology of Z-groups). Page 68 ends « TSVP ».
%
% What the batch is about. Page 61 finishes the proof begun at the foot of
% page 60 (batch 3): X' ⊂ X divisorial and Y-transversal at x iff X' is
% Y-flat at x and X' ⊗ k(y) divisorial (necessity, then sufficiency via
% Nakayama), then « Cor. Stabilité par extension de la base ». Pages 62-66,
% « Groupe de Brauer global » (complex analytic setting): the sequences
% 1 -> C^* -> Gl(n) -> GP(n-1) -> 1 and 1 -> Z/nZ -> Sl(n) -> GP(n-1) -> 1,
% the obstruction ∂(γ) ∈ H^2(X, Z/nZ); the monoid of classes of algebra
% bundles, the Brauer group Br(X) and the injection Br(X) -> H^2(X, C^*) into
% torsion; Questions 1-2 (is every torsion class, every class of
% H^2(X, Z/nZ), of the form ∂a?) and a reduction via the universal covering to
% group cohomology of π = π_1. Pages 68-70, under the cover « Restriction des
% scalaires »: for k'/k finite and X_1 = ∏_{S'/S} X', H^0, H^1 (principal
% homogeneous bundles) and Galois H^i commute with restriction of scalars;
% for k'/k radicial, (G_m)_1 = R(G_m) with unipotent quotient U, H^1(k,U) ≃
% Z/pZ over k_0((t)), the kernel V of the norm-like map not smooth,
% V ≃ ∏ μ_p when k'^p ⊂ k, and « (Z_1)_réd n'est pas définissable sur k ».
% Pages 72-80, under the cover « Petits fourbis cohomologiques »: H^1(Γ, A^{L*})
% = 1 for a finite-dimensional algebra A (twisted Hilbert 90, via a
% Zariski-density argument for k infinite; the finite case begun);
% Proposition (Lang): H^1(Γ, G^L) = {e} for G connected over a finite field,
% also for twisted forms, reduced to Lemme 1 (x ↦ x(ψ x^{(q)})^{-1} is
% surjective); the cohomology H^1(T, G) = colim H^1(Z_n, G_n) of a group with
% an automorphism T, Lemmes 1-5 (injectivity and surjectivity of
% H^1(T,G) -> H^1(T,G'')); page 80, twisting a Γ-set M by a 1-cocycle f,
% γ_f m = f(γ).γm, and the twisted Γ-group structure on G.
%
% Notation fixed once for the batch, following batches 2-3: Gl(n), Sl(n),
% GP(n-1) as written (GP = projective group); C^* underlined where he
% underlines it (sheaf of germs), set \underline{\mathbb{C}}^{*}; his
% double-struck B for the monoid of classes is \mathbb{B}(X); μ_p is
% \mu_p; the restriction of scalars X_1 = \prod_{S'/S} X'; his γ with a
% small arrow above and f below is set \gamma_f.
%
% Readings to check first: page 61 (fast); the margins of pages 65, 68, 69,
% 70 and 73; the struck blocks of pages 73-75; the interlinear additions of
% pages 76 and 78.
% Apparatus: 211 \ill{}, 29 \uncertain{}.

\documentclass[11pt,a4paper]{article}
% Le préambule commun à toutes les transcriptions du fonds Grothendieck.
%
% Il tient en une idée : **l'incertitude se compose, elle ne se résout pas.**
% Une transcription qui lisse un mot douteux a détruit la seule chose pour
% laquelle on la faisait — un lecteur doit pouvoir savoir, à chaque endroit,
% ce qui était sur le feuillet et ce qui a été ajouté. Les six macros qui
% suivent sont donc l'appareil critique, et non de la décoration.
%
% Les mêmes commandes sont reconnues par `scripts/render.mjs`, qui produit la
% vue de lecture du site. Une macro ajoutée ici doit l'être aussi là-bas,
% faute de quoi le rendu échouera bruyamment — ce qui est voulu : un rendu
% silencieusement faux est pire qu'un rendu absent.

\ProvidesPackage{grothendieck}[2026/08/08 Transcriptions du fonds Grothendieck]

\usepackage{amsmath,amssymb,amsthm}
\usepackage{mathrsfs}
\usepackage{stmaryrd}
% Les diagrammes commutatifs sont partout dans ce fonds : c'est la langue
% même de la géométrie algébrique de Grothendieck, et une transcription qui
% les rendrait en prose ne transcrirait rien.
\usepackage{tikz-cd}
% Français par défaut — c'est la langue des feuillets — mais l'anglais est
% chargé aussi : la traduction et le résumé sont des documents anglais, et la
% typographie française (espaces avant les deux-points) y serait une faute.
% Ces documents basculent par \selectlanguage{english} en tête de corps.
\usepackage[english,french]{babel}
\usepackage{fontspec}
\usepackage{geometry}
\geometry{a4paper,margin=25mm,marginparwidth=28mm}
\usepackage{marginnote}
\usepackage{xcolor}
% ulem, not soul. `soul`'s \st and \ul are built on a character-by-character
% scan that XeLaTeX's font machinery breaks: the document fails with an
% undefined control sequence rather than degrading. ulem does the same two jobs
% and is Unicode-engine safe. `normalem` stops it hijacking \emph, which the
% transcriptions use for Grothendieck's own emphasis.
\usepackage[normalem]{ulem}
\usepackage{hyperref}
\hypersetup{colorlinks=true,linkcolor=black,urlcolor=black,pdfborder={0 0 0}}

% --- Métadonnées du lot -----------------------------------------------------
% Reprises telles quelles par le script de rendu, qui les affiche en tête de
% la vue de lecture. Elles fixent ce que le fichier prétend couvrir : sans
% elles, une transcription flotte sans point d'attache dans un fonds de seize
% mille feuillets.

\newcommand{\folder}[1]{\gdef\grothFolder{#1}}
\newcommand{\batch}[1]{\gdef\grothBatch{#1}}
\newcommand{\leaves}[2]{\gdef\grothFirst{#1}\gdef\grothLast{#2}}
\newcommand{\foldertitle}[1]{\gdef\grothFoldertitle{#1}}
\gdef\grothFolder{?}\gdef\grothBatch{?}\gdef\grothFirst{?}\gdef\grothLast{?}\gdef\grothFoldertitle{}

% --- Le résumé --------------------------------------------------------------
%
% Il ouvre la lecture modernisée, sous le titre « Résumé », et il est composé
% en petit. Ce n'est pas un ornement : c'est le seul passage du document qui
% ne soit pas des mathématiques, et un lecteur doit pouvoir le reconnaître
% avant de l'avoir lu — puis le sauter s'il connaît déjà le sujet, ou s'y
% arrêter s'il ne le connaît pas. Le filet qui le suit marque la reprise du
% corps.

\newenvironment{resume}
  {\small\setlength{\parskip}{0.45em}}
  {\par\medskip\hrule\medskip}

% --- Mots-clés ---------------------------------------------------------------
%
% En anglais, en clôture du résumé de la lecture modernisée : le vocabulaire
% moderne sous lequel chercher ce que les feuillets construisent. Le manifeste
% du site (`npm run manifest`) les extrait pour étiqueter la cote entière —
% cette ligne est la source unique des tags, il n'y a pas de fichier à côté.
\newcommand{\keywords}[1]{\par\smallskip\noindent
  {\footnotesize\textsc{Keywords} --- \textit{#1}}\par}

% --- L'appareil critique ----------------------------------------------------

% Le numéro de feuillet, en marge. C'est l'ancre qui permet de régler un
% désaccord sur une formule en regardant *un* feuillet plutôt qu'une chemise
% de six cents. Le site s'en sert aussi pour tourner les pages du fac-similé
% quand on fait défiler la transcription.
\newcommand{\leaf}[1]{%
  \marginnote{\footnotesize\color{gray!70}\textsf{#1}}%
  \label{leaf:#1}\ignorespaces}

% Illisible : on ne devine pas. Un mot inventé qui se lit comme les autres est
% le pire résultat possible.
\newcommand{\ill}{\textcolor{red!60!black}{[\,\ldots\,]}}

% Lecture douteuse : le mot est donné, et signalé comme incertain.
\newcommand{\uncertain}[1]{\uline{#1}}

% Ajout de l'éditeur — une lettre manquante, un mot sous-entendu.
\newcommand{\add}[1]{\textcolor{blue!50!black}{[#1]}}

% Biffé par Grothendieck lui-même. Ce qu'il a rayé fait partie du document :
% une rature dit souvent où la pensée a changé de direction.
\newcommand{\struck}[1]{\sout{#1}}

% Note du transcripteur, sur l'état matériel du feuillet.
\newcommand{\note}[1]{\par\smallskip{\small\color{gray!60!black}\textit{[#1]}}\par\smallskip}

% Note marginale de Grothendieck, distincte du corps du texte.
\newcommand{\marginal}[1]{\par\smallskip\noindent\hspace{1em}%
  {\small\color{gray!60!black}#1}\par\smallskip}

% --- Titre ------------------------------------------------------------------

\AtBeginDocument{%
  \begin{center}
    {\large\bfseries Fonds Alexandre Grothendieck --- cote n\textsuperscript{o}~\grothFolder}\\[2pt]
    {\itshape\grothFoldertitle}\\[4pt]
    {\small feuillets \grothFirst--\grothLast\ (lot \grothBatch)}\\[2pt]
    {\footnotesize\color{gray!60!black}Transcription établie d'après le fac-similé numérisé
     par l'Université de Montpellier.\\
     Les lectures incertaines sont soulignées, les illisibles notées [\,\ldots\,],
     les ajouts entre crochets.}
  \end{center}
  \vspace{1em}}

\endinput


\folder{46}
\batch{4}
\pages{61}{80}
\dating{[vers 1966-vers 1972]}
\watermark{Édition de démonstration}
\foldertitle{Schémas en groupes et fibrés principau[x] (Divers) : notes manuscrites (s.d.), tapuscrit (s.d.)}

\begin{document}

\page{61}

\note{suite de la démonstration commencée au bas de la page 60}
est exacte, donc, comme $\mathcal{O}_X$ est plat sur $Y$, il s'ensuit
(\ill{}) que $\mathcal{O}_{X'}$ est plat sur $Y$, et que
$\mathcal{O}_{X'} \otimes k(z)$ est \add{\uncertain{un quotient}} divisoriel
par $\mathcal{O}_X \otimes k(z)$.

\emph{\uncertain{Suffisance}}\note{le mot est encadré ; au-dessus, « Faisons
\ill{} $z = y$ »}
\quad Suite exacte $0 \to \underline{I} \to \mathcal{O}_X \to
\mathcal{O}_{X'} \to 0$, d'où, puisque $\mathcal{O}_{X'}$ est $Y$-plat
\emph{en $x$}, la suite exacte
\[
  0 \to \underline{I} \otimes k(y) \to \mathcal{O}_X \otimes k(y) \to
  \mathcal{O}_{X'} \otimes k(y) \to 0 ,
\]
d'autre part $\mathcal{O}_{X'} \otimes k(y)$ est divisoriel \add{en $x$}, il
s'ensuit que $(\underline{I} \otimes k(y))_x$ est monogène, donc (Nakayama)
$\underline{I}$ \struck{\ill{}} est \struck{\ill{} \ill{} \ill{}}
\add{monogène} au voisinage de $x$, soit $\underline{I} = \varphi
\mathcal{O}_X$ [on rapetisse $X$]. Comme $\varphi$ dans
$\mathcal{O}_X \otimes k(y)$ est un générateur de $\underline{I} \otimes k(y)$,
il n'est pas diviseur de $0$ dans $\mathcal{O}_X \otimes k(y)$, i.e.\ la suite
\[
  0 \to \mathcal{O}_X \otimes k(y) \xrightarrow{\ \varphi\ }
  \mathcal{O}_X \otimes k(y) \to \mathcal{O}_{X'} \otimes k(y) \to 0
\]
\ill{} \ill{} la suite
$0 \to \mathcal{O}_X \xrightarrow{\ \varphi\ } \mathcal{O}_X \to
\mathcal{O}_{X'} \to 0$ est \emph{exacte}. Je dis que
$\varphi : \mathcal{O}_X \to \mathcal{O}_X$ est \emph{injectif} en $x$ : soit
$I$ le noyau, on aura alors les suites exactes
$0 \to I \to \mathcal{O}_X \xrightarrow{\ \varphi\ } \ill{} \to 0$,
$0 \to \underline{I} \to \mathcal{O}_X \to \mathcal{O}_{X'} \to 0$, où
$\underline{I}$ est plat (car $\mathcal{O}_X$ et $\mathcal{O}_{X'}$ le sont)
donc on a les suites exactes
$0 \to I \otimes k \to \mathcal{O}_X \otimes k \to \underline{I} \otimes k \to
0$, d'où $I \otimes k = 0$ d'où $I = 0$ (Nakayama).
\note{le terme illisible de la première suite est couvert d'une tache
d'encre}

\ill{} \ill{} \emph{Cor} \quad Stabilité par extension de la base.

\page{62}

\emph{Groupe de Brauer global.}

Pour avoir des fondements solides et bien connus, nous nous plaçons en
Géométrie \struck{Algébrique} \add{analytique} [sur le corps $\mathbb{C}$ des
complexes]. Les groupes de cohomologie seront pris p.r.\ à la topologie
usuelle.

Des \struck{\ill{}} \add{diagrammes de} suites exactes

\begin{tikzcd}[column sep=small, row sep=small]
e \arrow[r] & \mathbb{C}^{*} \arrow[r] & \mathrm{Gl}(n) \arrow[r] & \mathrm{GP}(n-1) \arrow[r] & e \\
e \arrow[r] & \mathbb{Z}/n\mathbb{Z} \arrow[r] \arrow[u, hook] & \mathrm{Sl}(n) \arrow[r] \arrow[u] & \mathrm{GP}(n-1) \arrow[r] \arrow[u, no head] & e
\end{tikzcd}

on tire (pour un bon $X$) la « suite exacte »

\begin{tikzcd}[column sep=small, row sep=small, nodes={font=\scriptsize}]
\cdots \arrow[r] & H^1(X, \underline{\mathbb{C}}^{*}) \arrow[r] & H^1(X, \underline{\mathrm{Gl}}(n)) \arrow[r] & H^1(X, \mathrm{GP}(n-1)) \arrow[r] & H^2(X, \underline{\mathbb{C}}^{*}) \\
\cdots \arrow[r] & H^1(X, \mathbb{Z}/n\mathbb{Z}) \arrow[r] \arrow[u] & H^1(X, \underline{\mathrm{Sl}}(n)) \arrow[r] \arrow[u] & H^1(X, \mathrm{GP}(n-1)) \arrow[r] \arrow[u, no head] & H^2(X, \mathbb{Z}/n\mathbb{Z}) \arrow[u]
\end{tikzcd}

\note{dans les deux diagrammes, la flèche verticale sous $\mathrm{GP}(n-1)$ est
un double trait (égalité)}

considérons, pour un fibré donné $\gamma \in H^1(X, \mathrm{GP}(n-1))$,
l'élément
\[
  \partial(\gamma) \in H^2(X, \mathbb{Z}/n\mathbb{Z})
\]
obstruction à remonter le groupe structural à $\underline{\mathrm{Sl}}(n)$.
Alors l'obstruction à remonter le groupe structural à $\mathrm{Gl}(n)$ est
l'image de $\partial(\gamma)$ dans $H^2(X, \mathbb{C}^{*})$ ; elle est donc
d'ordre $n$ [l'ensemble des éléments d'ordre $n$ de $H^2(X, \mathbb{C}^{*})$
étant exactement l'image de $H^2(X, \mathbb{Z}/n\mathbb{Z})$].

Les classes de fibrés sur $X$, de groupe $\mathrm{GP}(n-1)$, ne sont autres
que les classes de fibrés en algèbres \struck{\ill{}} de \struck{\ill{}}
\add{rang} $n^2$ ; celles qui proviennent de $\mathrm{Gl}(n)$ sont les fibrés
en algèbres \uncertain{qui}

\page{63}

sont de la forme $L(E, E)$, où $E$ est un fibré vectoriel de dim.\ $n$.

Si $\mathcal{A}$, $\mathcal{B}$ sont deux fibrés en algèbres de
\struck{\ill{}} \add{rang} $m^2$, $n^2$, alors $\mathcal{A} \otimes
\mathcal{B}$ est un fibré en algèbres de rang $m^2 n^2 = (mn)^2$, donc définit
un élément de $H^1(X, \mathrm{GP}(mn-1))$. Cela est associé à une
représentation de groupes
\[
  \mathrm{GP}(m-1) \times \mathrm{GP}(n-1) \longrightarrow \mathrm{GP}(mn-1)
\]
déduite de $\mathrm{Gl}(m) \times \mathrm{Gl}(n) \longrightarrow
\mathrm{Gl}(mn)$ par passage au quotient. Ainsi l'ens.
\[
  \mathbb{B}(X) = \coprod_{n \geqslant 1} H^1(X, \mathrm{GP}(n-1))
\]
est un \struck{\ill{}} monoïde \add{\uncertain{commut., associatif}}, dont
l'élément unité est l'unique élément de degré $1$. On \struck{\ill{}}
introduit dans ce monoïde une relation d'équivalence : deux \ill{} \ill{}
$a$, $b$ sont équivalents, si on peut trouver $c$, $c$ étant « trivial » ---
i.e.\ provenant d'un fibré vectoriel --- t.q.\ $ac = bc$. Si $a \sim b$ et
$a' \sim b'$, alors $aa' \sim bb'$, donc le monoïde quotient, noté
$\mathrm{Br}(X)$, est un monoïde \struck{\ill{}} (ass., comm., avec $1$).
C'est même un groupe, car si
\note{le « $n \geqslant 1$ » sous le coproduit est surchargé : le « $1$ » est
écrit sur un autre chiffre}

\page{64}

$a \in \mathbb{B}(X)$, soit $a^{\circ}$ la classe « opposée » (en remplaçant
une algèbre par l'algèbre opposée) alors $aa^{\circ}$ est trivial. En effet,
on a canoniquement
\[
  \mathcal{A} \otimes \mathcal{A}^{\circ} \simeq L(\mathcal{A}, \mathcal{A}) :
\]
D'où la définition du \emph{groupe de Brauer}, noté $\mathrm{Br}(X)$.

On voit que
\[
  \delta(ab) = i_m\,\delta(a) + i_n\,\delta(b)
\]
\ill{} \ill{} \ill{} $\mathbb{Z}/n\mathbb{Z}$, $\mathbb{Z}/m\mathbb{Z}
\to \mathbb{Z}/mn\mathbb{Z}$, d'où résulte qu'on a un homomorphisme de
\struck{groupes} \add{monoïdes}
\[
  \mathbb{B}(X) \longrightarrow H^2(X, \mathbb{Q}/\mathbb{Z})
\]
\note{le $\mathbb{B}$ est entouré et surcharge un « Br »}
cet homomorphisme définit, par composition
$\mathbb{B}(X) \to H^2(X, \underline{\mathbb{C}}^{*})$, d'où \ill{}
\struck{\ill{} \ill{}} \ill{} \struck{\ill{} \ill{} \ill{}}
\add{formé des} classes « triviales ». \struck{\ill{}} D'où un homom.\
naturel \emph{injectif}
\[
  \boxed{\mathrm{Br}(X) \xrightarrow{\ \varphi\ } H^2(X, \underline{\mathbb{C}}^{*})}
\]
dont l'image est contenue dans $\operatorname{Im} H^2(X,
\mathbb{Q}/\mathbb{Z})$, donc dans le sous-groupe de Torsion.

\emph{Question 1} \quad Tout élément de torsion de $H^2(X,
\underline{\mathbb{C}}^{*})$ provient-il d'un élément de $\mathrm{Br}(X)$ ?
\ill{} \ill{} \ill{} \ill{} ? \uncertain{question} \uncertain{suivante} plus
précise :

\emph{Question 2} \quad Tout élément de $H^2(X, \mathbb{Z}/n\mathbb{Z})$
est-il de la forme $\partial a$, avec $a \in H^1(X, \mathrm{GP}(n-1))$ ???

\page{65}

[Il faudrait d'abord résoudre cette question du point de vue
\emph{topologique}. Ce qui revient d'abord au cas de la classe canonique dans
$H^2(Z, \mathbb{Z}/n\mathbb{Z}, \mathbb{Z}/n\mathbb{Z})$ : est-elle de la
forme $\delta(a)$, avec $a \in H^1(Z, \mathbb{Z}/n\mathbb{Z},
\mathrm{GP}(n-1))$ ?
\note{le premier argument est une lettre isolée, lue $Z$, avec un indice ou
une virgule peu nets ; ce qui suit le premier « $H^1($ » est surchargé}
\marginal{\uncertain{faisceau} des \uncertain{sections} à \ill{}}
\note{note cerclée dans la marge droite, reliée par une flèche au
$\mathrm{GP}(n-1)$ souligné}
Sinon, on aurait \ill{} \ill{} \ill{} \ill{} à : $H^2(X,
\mathbb{Q}/\mathbb{Z})$, tout élément est-il l'obstruction d'un élément de
$H^1(X, \underline{\mathrm{GP}(n-1)})$ pour $n$ convenable ??

Puis on \ill{} \ill{} \ill{} \ill{} se ramener à voir si les éléments de
torsion de $H^3(X, \mathbb{Z})$ \struck{$\simeq H^2(X, \ill{})$} [i.e.\ les
éléments de $\partial(H^2(X, \mathbb{Q}/\mathbb{Z}))$] est l'obstruction
d'un élément de $H^1(X, \mathrm{GP}(n-1))$, avec $n$ convenable [N.B.
$H^3(X, \mathbb{Z}) \simeq H^2(X, \underline{\mathbb{C}}^{*})$] --- c'est la
question 1 \emph{topologique} ---]

On peut introduire le revêtement universel $\widetilde{X}$ de $X$, de groupe
de Galois $\pi$, et on a alors
\[
  0 \to H^2(\pi, \mathbb{Z}/n\mathbb{Z}) \to H^2(X, \mathbb{Z}/n\mathbb{Z})
  \to H^2(\widetilde{X}, \mathbb{Z}/n\mathbb{Z})^{\pi} \to \cdots
\]
et on peut essayer de montrer que du moins les éléments de $H^2(X,
\mathbb{Z}/n\mathbb{Z})$ qui proviennent de $H^2(\pi, \mathbb{Z}/n\mathbb{Z})$
sont dans $\operatorname{Im} H^1(X, \mathrm{GP}(n-1))$ ; \struck{ainsi donc}
il suffit pour cela de prouver que $H^1(\pi, \mathrm{GP}(n-1)) \to H^2(\pi,
\mathbb{Z}/n\mathbb{Z})$ est surjectif ---
\struck{\ill{} \ill{} \ill{} prenant $X = K(1, \pi)$ \ill{}}
Du moins, est-il vrai que tout élément de $H^2(\pi, \mathbb{Q}/\mathbb{Z})$
est dans $\operatorname{Im} H^1(\pi, \mathrm{GP}(n-1))$ pour $n$
\emph{convenable} ?
\note{après « surjectif », une ligne et demie est annulée par un gribouillis
en boucles ; on y lit seulement « $X = K(1, \pi)$ »}

\page{66}

Et si ceci aussi est faux, est-il donc vrai au moins que tout élément de
$H^2(\pi, \mathbb{C}^{*})$ est dans $\operatorname{Im} \struck{\ill{}} H^1(\pi,
\mathrm{GP}(n-1))$ pour $n$ convenable ??

---\,]
\note{le crochet fermant, isolé en bas à droite de la page, clôt celui qui
ouvre la page 65}

\section*{« Restriction des scalaires » dans les groupes algébriques et leurs espaces principaux homogènes\ldots}

\note{inscrit et souligné, avec les guillemets, en haut à droite d'une chemise
(page 67), qui ne reçoit pas de numéro de page}

\page{68}

$k$ corps, $k'$ extension finie de $k$, $S = \operatorname{Spec} k$,
$S' = \operatorname{Spec} k'$

\marginal{N.B. \ill{} \ill{} \ill{} \ill{} : voir dans la question\ldots}

a) $X'$ préschéma sur $k'$, $X_1 = \prod_{S'/S} X'/S'$

(i) \quad $\boxed{H^0(k, X_1) \simeq H^0(k', X')}$

(ii) \quad Si $X'$ est un schéma en groupes, \uncertain{noté} $G'$,
\emph{\ill{} sur $k$}, on a

\quad $\boxed{H^1(k, G_1) \simeq H^1(k', G')}$ \quad (classes de fibrés
principaux homogènes ----)

(iii) \quad Si $G'$ est un schéma en groupe \emph{commut.}, noté $M'$, on a

\quad $\boxed{H^i(k, M_1) \simeq H^i(k', M')}$ \quad pour tout $i$
\quad (Cohomologie galoisienne\ldots)

(i) Est vrai par définition (on n'a pas besoin de corps),

(ii) Soit $P'$ fibré principal homogène sous $G'$, alors
$P_1 = \prod_{S'/S} P'$ est un schéma sur $k$, sur lequel $G'$ opère. On voit
que $P_1$ est formellement fibré princ.\ homog.\ sous $G_1$, d'autre part
\struck{il est $\neq \emptyset$} pour utiliser la \uncertain{simplicité} de
$G'/k'$, on voit que $P_1 \neq \emptyset$, \struck{donc c'est un fibré
principal homogène}. La fonction $P' \mapsto P_1$ est pleinement fidèle.
\add{i.e.\ \ill{} \ill{} $P_1$ \ill{} \ill{} \ill{} catégories}
On montre que c'est une équivalence de catégories, un argument de descente
montre \ill{} \ill{} \ill{} \ill{} $P_1$ est trivial ----
\marginal{\uncertain{Contre-exemple} \ill{} si $G'$ \ill{} \ill{} \ill{} ??
\ill{} \ill{} $G'$ \ill{} \ill{} $k'$ ??}
\note{la note marginale, en quatre lignes obliques, est écrite en travers de
la marge gauche en regard de (ii) ; l'addition « i.e.\ldots\ catégories » est
entourée}

(iii) Soit $T'/S$ \struck{fini étale} galoisien, \ill{} \ill{}.
\[
  C^{*}(T/S, M_1) \simeq C^{*}(T'/S', M')
\]
d'où
\[
  H^{*}(T/S, M_1) \simeq H^{*}(T'/S', M') ,
\]
d'où, si $T$ parcourt certains $S$-préschémas,
\[
  \varinjlim_{T} H^{*}(T/S, M_1) \simeq \varinjlim_{T} H^{*}(T'/S', M') ,
\]
Or, si $T$ parcourt les rev.\ étales \uncertain{complets} (ou les produits
\ill{} \ill{} \uncertain{finis} sur $S$) les $T'/S'$ sont \emph{cofinaux} dans
les \add{\uncertain{catégories}} \struck{\ill{}} \ill{} rel.\ à $S'$ ! D'où la
conclusion.
\note{au bas de la page, à gauche : « TSVP »}

\page{69}

Supposons que $k'/k$ est \emph{radiciel}, prenons $M' = \mathbb{G}_{m,k'}$,
alors on a une suite exacte
\[
  e \to \mathbb{G}_m \to (\mathbb{G}_m)_1 \to U \to e ,
  \qquad (\mathbb{G}_m)_1 = \underline{\mathrm{Hom}}_S(S', \mathbb{G}_m)
\]
\note{l'égalité est écrite verticalement sous $(\mathbb{G}_m)_1$ ; la lecture
« $\underline{\mathrm{Hom}}$ » est douteuse}
où $U$ est unipotent connexe, de \ill{} \ill{} sur $k$ [en car.\ $2$, on peut
trouver $U$ de dim.\ $1$, \struck{\ill{}} par ex.\ si $[k':k] = 2$]. La suite
exacte de cohomologie galoisienne sur $k$ donne alors, compte tenu de
$H^i(k, (\mathbb{G}_m)_1) \simeq H^i(k', \mathbb{G}_m)$, l'isom.
\[
  H^1(k, U) \simeq \operatorname{Ker}\bigl(H^2(k, \mathbb{G}_m) \to
  H^2(k', \mathbb{G}_m)\bigr)
\]
Prenons p.\ ex.\ $k = k_0((t))$, $k' = k_0((t^{1/p}))$, où $k_0$ est un corps
fini de car.\ $p$ \ill{} \ill{}, on trouve par la \ill{} du \uncertain{corps
de classes local}
\[
  H^1(k, U) \simeq \mathbb{Z}/p\mathbb{Z} \neq 0
\]
\note{les parenthèses doubles de $k_0((t))$ se lisent aussi bien comme des
crochets}
[bien que, pour $p = 2$, $U$ soit une forme de $\mathbb{G}_a$ ; bien sûr,
cette forme n'est pas isomorphe à $\mathbb{G}_a$].
\marginal{N.B. On \uncertain{sait} \ill{} que $pU = 0$ (\ill{} \ill{} \ill{}
par $k'^{p} \subset k$)}

Notons aussi que \emph{l'ext.\ de $U$ par $\mathbb{G}_m$} [qui splitte sur
$\bar{k}$] \emph{n'est pas triviale}, car autrement on aurait
\[
  H^1(k, (\mathbb{G}_m)_1) \simeq H^1(k, \mathbb{G}_m) \times H^1(k, U)
\]
d'où $H^1(k, U) = 0$.
\note{sous $H^1(k, (\mathbb{G}_m)_1)$ et $H^1(k, \mathbb{G}_m)$, il note
« $= 0$ » verticalement}
Donc « la \ill{} \add{\uncertain{multiplicative}} de $(\mathbb{G}_m)_1$
\emph{n'est pas définie sur $k$} »

\page{70}

Regardons maintenant l'\emph{hom.\ \ill{}} et soit $V$ son noyau
\[
  e \to V \to (\mathbb{G}_m)_1 \to \mathbb{G}_m \to e
\]
\emph{Mais} $V$ \emph{n'est pas lisse sur $k$}, \struck{\ill{} \ill{}}
\marginal{dans le cas particulier envisagé où $H^1(k, U) \neq 0$}
\struck{Mais [dans le cas où \ill{} \ill{}]}
\ill{} \ill{} que $V_{\mathrm{réd}}$ [est exactement le radical
\add{unipotent} de $\mathbb{G}_m$]. D'ailleurs, dans le cas où $[k':k] = p$,
on voit facilement que $V$ est réduit. D'ailleurs, dans le cas où
$k'^{p} \subset k$, on voit facilement que \emph{$V$ est aussi le noyau de la
puissance $p$-ième dans $\mathbb{G}_m$}, donc on a un isom.
\[
  V \simeq \prod_{S'/S} (\mu_p)_{S'} .
\]
\marginal{\uncertain{précisions} \ill{} \ill{}\ldots}
\note{le paragraphe « D'ailleurs\ldots\ isom. » est marqué d'un double trait
vertical à gauche, la note marginale écrite de biais}
\struck{En tous cas}, on \ill{} \ill{} \ill{} \ill{} \ill{} une suite exacte
\[
  0 \to \mu_p \to V \to U \to 0
\]
(extension non triviale ---). Notons maintenant que si $G'$ est un groupe
alg.\ sur $k'$ \add{\ill{} \ill{}} \struck{\ill{}} ayant un centre isomorphe à
$(\mu_p)_{k'}$ [p.\ ex.\ $\mathrm{Sl}(p)_{k'}$] alors le centre de $G_1$ est
\[
  Z_1 = V
\]
donc $(Z_1)_{\mathrm{réd}}$ \emph{n'est pas définissable sur $k$}.

\section*{Petits fourbis cohomologiques}

\note{inscrit et souligné en haut à droite d'une chemise (page 71), qui ne
reçoit pas de numéro de page ; la lecture « fourbis » est douteuse. Les mêmes
mots sont répétés au crayon, de biais, dans la marge gauche de la page 72}

\page{72}

\emph{Proposition} \quad Soit $k$ un corps commutatif, $L$ une extension
galoisienne de $k$ de groupe de Galois $\Gamma$, $\underline{A}$ une algèbre
avec $1$ de dimension finie sur $k$, $\underline{A}^{L}$ l'algèbre sur $L$
obtenue par extension des scalaires à $L$, $\underline{A}^{L}$ est considérée
comme munie du groupe d'automorphismes défini par $\Gamma$. Soit
$G = \underline{A}^{L*}$ le groupe des éléments inversibles de
$\underline{A}^{L}$, $G$ \struck{$\underline{A}^{L}$} un $\Gamma$-groupe.
$H^1(\Gamma, G)$ est réduit à l'élément unité, ainsi que $H^1(\Gamma,
G^{\psi})$ si $\psi$ est un hom.\ croisé de $\Gamma$ dans \struck{\ill{} $G$}
\add{$\mathrm{Aut}\,\underline{A}$}. \emph{\ill{}} \quad $H^1(\Gamma,
\mathrm{Gl}(n, L))$ est réduit à l'élément unité.
\note{l'énoncé est marqué d'un trait vertical à gauche}

\emph{Démonstration} \quad $\Gamma$ opère dans $G^{\psi} = G$ par
$\sigma(\gamma)g = \psi(\gamma).g^{\gamma}$. Soit $\varphi \in Z^1(\Gamma,
G^{\psi})$ \ill{} \ill{}
\[
  x = \sum_{\gamma \in \Gamma} \varphi(\gamma)\,
  \sigma(\gamma).A \qquad (A \in \underline{A}^{L})
\]
\struck{\ill{} \ill{} \ill{}, $\sigma(\gamma)$ \ill{} \ill{} \ill{} \ill{}}
\note{dans cette formule et la suivante, une parenthèse biffée et illisible
suit le signe somme}
\[
  \sigma(\gamma)x = \sum_{\gamma' \in \Gamma}
    \sigma(\gamma)\varphi(\gamma')\,\sigma(\gamma\gamma').A
\]
\struck{$= \sum_{\gamma' \in \Gamma} (\sigma(\gamma')A)\,
\sigma(\gamma)\varphi(\gamma^{-1}\gamma')$}
\begin{align*}
  &= \sum_{\gamma' \in \Gamma} \sigma(\gamma)\varphi(\gamma^{-1}\gamma')\,
    \sigma(\gamma').A \\
  &= \sigma(\gamma)\varphi(\gamma^{-1}) \sum_{\gamma' \in \Gamma}
    \varphi(\gamma')\,\sigma(\gamma').A = \sigma(\gamma)\varphi(\gamma^{-1})x
\end{align*}
\[
  x = \varphi(\gamma^{-1})\,\sigma(\gamma)^{-1}x \qquad
  x = \varphi(\gamma)\,\sigma(\gamma)x
\]
\struck{$\varphi(\gamma) =$} D'où, si \emph{$x$ inversible}, la formule
\[
  \varphi(\gamma) = x\,(\sigma(\gamma)x)^{-1}
\]
Il reste à prouver donc que l'on peut trouver $A$ de telle façon que
$\sum \varphi(\gamma)\,\sigma(\gamma).A$ soit

\page{73}

inversible. Or, $\sigma(\gamma)A = \psi(\gamma).A^{\gamma}$, où $\psi(\gamma)$
est une transformation linéaire de $\underline{A}^{L}$, si $x$ n'était pas
inversible, on aurait ($P$ étant la fonction \add{polynôme}
$P(A) = \struck{\det(\ill{})}$ \add{\uncertain{déterminant de}} $f_A$, \ill{}
de la transformation $B \mapsto AB$ \uncertain{ou} $B \mapsto BA$)
\[
  P\Bigl(\sum_{\gamma} \varphi(\gamma)\bigl(\psi(\gamma).A^{\gamma}\bigr)\Bigr) = 0
\]
\struck{\ill{} \ill{} \ill{} \ill{} \ill{} \ill{} $= 0$, i.e.\ \ill{}
$\sum_{\gamma}$ \ill{} \ill{} $A$.}
Considérons la fonction sur $(\underline{A}^{L})^{\Gamma}$ définie par
\struck{$\bar{P}((A_\gamma)) = P(\sum_{\gamma} \varphi(\gamma) \ldots$}
\[
  \bar{P}((A_\gamma)) = P\Bigl(\sum_{\gamma} \varphi(\gamma)\bigl(\psi(\gamma).
  A_{\gamma}^{\gamma}\bigr)\Bigr)
\]
Elle est nulle \ill{} \add{en faisant la substitution}
($A_\gamma = A^{\gamma}$ ($A \in \underline{A}^{L}$)) d'où on conclut qu'elle
est identiquement nulle (du moins si \emph{le corps $k$ est infini}).
\marginal{argument \ill{} \ill{} \ill{} \ill{} d'un corps !!!}
On aurait donc, en prenant $A_\gamma = 0$ si $\gamma \neq 1$, $A_1 = 1 \in
\underline{A}^{L}$ : $P(1) = 0$, ce qui est absurde.

Il reste : l'égalité \ill{} \ill{} \ill{} le cas \emph{fini}. Dans ce cas,
\struck{\ill{} \ill{} \ill{} \ill{} \ill{} fini ; $\mathbb{G}_m$ \ill{}
\ill{} \ill{} cyclique \ill{} \ill{} $\Gamma$} \ill{} \ill{}, et prouvons
$H^1(\Gamma, G^{\psi}) = e$ si $\psi$ est \ill{} $\varphi$ est \ill{} \ill{}
\ill{} \ill{} $\varphi(u) = a$
\begin{gather*}
  \varphi(u^0) = 1 \\
  \varphi(u) = a \\
  \varphi(u^2) = a\,\sigma(u)a \\
  \varphi(u^3) = a\,\sigma(u)a\,\sigma(u)^2 a \\
  \ldots \\
  \varphi(u^{n-1}) = a\,\sigma(u)a \cdots \sigma(u)^{n-2}a
\end{gather*}
\marginal{\ill{} \ill{} \ill{} \ill{} un hom.\ croisé de $\Gamma$ dans
$\mathrm{Aut}\,G^{L}$ \ill{} ; $\Gamma$ \uncertain{engendré par} $u$,
$u\lambda = \lambda^{q}$ ($q = [k]$) \ill{} \ill{}}
\note{la fin de la page est très raturée : deux lignes barrées, un mot
entouré, et la note marginale, de biais, reliée par une flèche au début de la
ligne « $\varphi$ est\ldots »}

\page{74}

\emph{Proposition} \quad Soit $G$ un groupe algébrique \emph{connexe} sur un
corps fini $k$, $L$ une extension finie de $k$, de groupe de Galois $\Gamma$,
$G^{L}$ le groupe obtenu à partir de $G$ par extension des scalaires à $L$ (il
est encore connexe sur $L$), $\Gamma$ opère dans $G^{L}$ comme groupe
d'automorphismes. \struck{Soit \ill{}}. Alors $H^1(\Gamma, G^{L}) = \lbrace
e \rbrace$. Plus généralement, soit $\mathrm{Aut}\,G^{L}$ le groupe des
automorphismes du groupe algébrique $G^{L}$ \emph{sur $L$}, et soit $\psi$ un
\struck{\ill{}} \add{hom.} croisé de $\Gamma$ dans $\mathrm{Aut}\,G^{L}$, soit
$G^{L\psi}$ le groupe $G^{L}$ muni de $\Gamma$ opérant par
$\sigma(\gamma).g = \psi(\gamma).g^{\gamma}$ ; alors, on a encore
\[
  H^1(\Gamma, G^{L\psi}) = \lbrace e \rbrace
\]
\marginal{à voir !}
\note{l'énoncé est marqué d'un trait vertical à gauche, la note marginale en
regard}

\emph{Démonstration} \quad Soit $q$ = nb d'éléments de $k$, et $n = [L:k]$.
$\Gamma$ est le groupe cyclique engendré par l'automorphisme
$f\lambda = \lambda^{q}$, on désigne par $u$ l'automorphisme correspondant de
$G^{L}$. Soit $\Omega$ une clôture algébrique de $L$, on désigne par
$\tilde{u}$ l'automorphisme \struck{\ill{}} de $G^{\Omega}$ défini \add{par
l'automorphisme $\lambda \mapsto \lambda^{q}$ de $\Omega$} (on a donc
$u^{n} = 1$, mais $\tilde{u}^{n} \neq 1$ ; l'ensemble des invariants de
$\tilde{u}^{n}$ dans $\Omega$ --- resp.\ $G^{\Omega}$ --- est $L$ --- resp.\
$G^{L}$ ---). On désigne par $\tilde{v}$ l'automorphisme
$x \mapsto \psi(f)x^{\tilde{u}}$ de $G^{\Omega}$, induisant un automorphisme
$v$ de $G$, qui n'est autre que $\sigma(f)$. On a, puisque $\psi(\gamma)$ est
un \struck{$\psi\,\psi^{\tilde{u}} \cdots \psi^{\tilde{u}}$} hom.\ croisé de
$\Gamma$ dans $\mathrm{Aut}\,G^{L}$
\note{la définition de $\tilde{u}$ a été reprise trois fois : on y voit,
biffés, « $f : \lambda \mapsto \lambda^{q}$ de $\Omega$ », « $x \mapsto
x^{\tilde{u}}$ » et, en surcharge, « $x \mapsto x^{q}$ »}

\page{75}

\[
  \psi^{(u)}\psi^{(u^2)} \cdots \psi^{(u^{n-1})} = 1
\]
(où $\psi^{(u^k)} = u^{k}\psi u^{k-1}$), d'autre part, en désignant par
$\tilde{\psi}$ le prolongement de $\psi$ à $G^{\Omega}$, et notant que les
$\psi^{(u^k)}$ sont encore des \struck{\ill{}} automorphismes \add{du groupe
algébrique} $G^{L}$ sur $L$, \struck{\ill{}} \add{\uncertain{induisant}} des
automorphismes $\tilde{\psi}^{(\tilde{u}^k)}$ de $G^{\Omega}$ :
\[
  (*) \qquad \tilde{\psi}^{\tilde{u}}\,\tilde{\psi}^{\tilde{u}^2} \cdots
  \tilde{\psi}^{\tilde{u}^{n-1}} = 1
\]
Soit $\varphi \in Z^1(\Gamma, G^{L\psi})$, donnée par $\varphi(u) = a \in
G^{L\psi}$, avec
\[
  (**) \qquad a\,a^{v}a^{v^2} \cdots a^{v^{n-1}} = 1
\]
On cherche $b \in G^{L}$ tel que $a = b(b^{v})^{-1}$ ; on va voir qu'il
suffit de \uncertain{trouver} $b \in G^{\Omega}$ tel que
$a = b(b^{\tilde{v}})^{-1}$ : on va voir qu'alors automatiquement
$b \in G^{L}$.

Pour ceci, on \struck{\ill{} \ill{} $a = b(b^{\tilde{v}})^{-1}$} :
\uncertain{écrivons} $(**)$
\note{suit un bloc annulé par un grand trait sinueux, où l'on lit
« $a^{v^k} = a^{\tilde{v}^k}$ »,
« $\tilde{v}^{k} = (\tilde{\psi}\tilde{u})^{k} = \tilde{\psi}
\tilde{\psi}^{\tilde{u}}\tilde{\psi}^{\tilde{u}^2} \cdots
\tilde{\psi}^{\tilde{u}^{k-1}}\tilde{u}^{k}$ (évidemment \ldots) » et
« $a^{v^k} = a^{\tilde{v}^k} = (\tilde{\psi} \cdots
\tilde{\psi}^{\tilde{u}^{k-1}})\,b\,b^{v} \ldots b\,b^{\tilde{v}^{-1}}$ »}
on constate que
\[
  a^{v^k} = a^{\tilde{v}^k} = b^{\tilde{v}^k}\bigl(b^{\tilde{v}^{k+1}}\bigr)^{-1}
\]
d'où on tire
\[
  b\bigl(b^{\tilde{v}^n}\bigr)^{-1} = 1 \qquad \text{soit} \qquad
  b^{\tilde{v}^n} = b
\]
Or on a
$\tilde{v}^{k} = (\tilde{\psi}\tilde{u})^{k} = \tilde{\psi}
\tilde{\psi}^{\tilde{u}} \cdots \tilde{\psi}^{\tilde{u}^{k-1}}\tilde{u}^{k}$
(récurrence sur $k$) d'où
\[
  \tilde{v}^{n} = \bigl(\tilde{\psi}\tilde{\psi}^{\tilde{u}} \cdots
  \tilde{\psi}^{\tilde{u}^{n-1}}\bigr)\tilde{u}^{n} = \tilde{u}^{n}
\]
en vertu de $(*)$, on a donc $b^{\tilde{u}^n} = b$, i.e.\ $b \in G^{L}$.

On est maintenant ramené au théorème \uncertain{géométrique} suivant :

\page{76}

\emph{Lemme 1} \quad Soit $G$ un groupe algébrique connexe sur le corps
algébriquement clos $\Omega$ de car.\ $p \neq 0$, \add{« défini » sur le corps
$\mathbb{F}_q$ à $q$ éléments, soit} \struck{soit} $q$ une puissance
\struck{\ill{}} $\psi$ un \add{\uncertain{automorphisme du groupe
algébrique}} \struck{\ill{}}, considérons \struck{l'automorphisme $\sigma$ du
groupe $G$ \add{défini} défini par} l'application de $G$ dans $G$ donnée par
\[
  x \longmapsto x\,\bigl(\psi x^{(q)}\bigr)^{-1}
\]
où $x^{(q)}$ est le transformé de $x$ par l'automorphisme qui correspond à
$\lambda \mapsto \lambda^{(q)}$ dans $\Omega$ (dont le corps des invariants
est $\mathbb{F}_q$) ; cette application est alors \emph{surjective}.
\note{l'énoncé est marqué d'un trait vertical à gauche}

\struck{On \ill{} \ill{} l'\ill{} \ill{} \ill{}}
Considérons la différentielle de cette application à l'élément neutre : par
un calcul presque formel, on trouve que c'est l'identité (car la
différentielle de $x \mapsto \psi x^{(q)}$ est la \struck{\ill{}}
\add{composée} de $d\psi(e)$ avec la différentielle de $x^{(q)}$, laquelle
est nulle).

Donc l'application envisagée a pour image un \emph{voisinage} de $e$. C'est
d'ailleurs un hom.\ croisé (quand $G$ opère sur lui-même par automorphismes
intérieurs par $\psi x^{(q)}$ : $\theta(xy) = \theta(x)\,
\mathrm{int}(\psi x^{(q)})\,\theta(y)$).

Si $G$ est abélien, l'image $\theta(G)$ est donc un sous-groupe, qui est
\add{\ill{} \ill{}} \struck{\ill{}} : $G$ puisque $G$ est connexe.

Dans le cas général, on \ill{}

\page{77}

Soit $G$ un groupe où opère $\mathbb{Z}$, \add{$T^{1} = T$}, \ill{}
\struck{\ill{}} \ill{}, soit $G_n$ l'ens.\ des éléments de $G$ tels que
$T^{n}x = x$. On suppose que $G = \bigcup G_n$ (les $G_n$ forment une suite
croissante). Alors $G_n$ est un $\mathbb{Z}_n$-groupe, et on peut considérer
sa cohomologie \add{\ill{}} \ill{}, pour $G$ commutatif, \ill{} \ill{} :
\begin{align*}
  H^0(\mathbb{Z}_n, G_n) &= G_1 \quad \text{ens.\ des éléments de $G$
    invariants par $\mathbb{Z}$} \\
  H^1(\mathbb{Z}_n, G_n) &= \text{ens.\ des éléments $\alpha$ de $G_n$ tels
    que} \\
  &\qquad \alpha\,T\alpha\,T^{2}\alpha \cdots T^{n-1}\alpha = 1 ,
\end{align*}
mod la relation d'équivalence définie par $\beta \simeq \alpha$ si et
seulement si il existe $u \in G_n$ tel que $\beta = u^{-1}\alpha\,Tu$.

D'ailleurs supposons que $n \mid m$, alors $\mathbb{Z}_n$ est un quotient de
$\mathbb{Z}_m$, \struck{\ill{}} par un sous-groupe $N_{n,m}$, et on a
$\boxed{G_n = G_m^{N_{n,m}}}$. On a donc fonctoriellement une suite exacte
\[
  e \to H^1(\mathbb{Z}_n, G_n) \xrightarrow{\ \mathrm{inf}\ }
  H^1(\mathbb{Z}_m, G_m) \to H^1(N_{n,m}, G_m)^{\mathbb{Z}_m} ,
  \qquad N_{n,m} \simeq \mathbb{Z}_{m/n}
\]
[Si $G$ est abélien, on a aussi \ill{} \ill{} $\mathbb{Z}_{m/n}$ \ill{}
$H^i(\mathbb{Z}_n, G_n)$, et on a des \add{homom.}
\[
  H^i(\mathbb{Z}_n, G_n) \to H^i(\mathbb{Z}_m, G_m)
\]
d'ailleurs, on sait à partir de $i = 1$, une périodicité de période $2$ dans
les $H^i(\mathbb{Z}_n, G_n)$, \struck{\ill{}} donc $i > 1$ n'importe pas grand
chose].

\emph{Lemme 1} \quad Soit $\alpha \in G$ tel que $\alpha\,T\alpha\,T^{2}\alpha
\cdots T^{n-1}\alpha = 1$, alors $\alpha \in G_n$. Supposons $\alpha, \beta
\in G$ satisfaisant aux conditions précédentes et $\beta = u^{-1}\alpha\,Tu$
($u \in G$), alors $u \in G_n$.
\note{l'énoncé est marqué d'un trait vertical à gauche}

En effet, dans le premier cas, appliquons $T$ :
$T\alpha\,T^{2}\alpha \cdots T^{n-1}\alpha\,T^{n}\alpha = 1$ i.e.\
$T^{n}\alpha$ est \emph{un inverse} de $T\alpha\,T^{2}\alpha \cdots
T^{n-1}\alpha$, comme \ill{} \ill{} \ill{} \ill{} $\alpha$, \ill{}
$T^{n}\alpha = \alpha$. Dans le second, écrivons $\beta\,T\beta \cdots
T^{n-1}\beta = 1$, et y remplaçons $\beta$ par $u^{-1}\alpha\,Tu$, on trouve

\page{78}

\note{en tête de page, deux lignes de sa main, barrées ensemble d'un long
trait ondulé : « F. Peterson \quad Functional Cohomology operations. Some
non-embedding Problems. » et « Moore \quad The double suspension and
$p$-primary components of homotopy groups of spheres »}
\[
  u^{-1}\alpha\,Tu\;Tu^{-1}\,T\alpha\,T^{2}u\;T^{2}u^{-1}\,T^{2}\alpha\,T^{3}u
  \cdots T^{n-1}u^{-1}\,T^{n-1}\alpha\,T^{n}u = 1
\]
\note{les facteurs qui se simplifient sont barrés de petits traits obliques}
i.e.\ $u^{-1}(\alpha\,T\alpha \cdots T^{n-1}\alpha)\,T^{n}u = 1$ d'où, comme
$(\ \ ) = 1$, $u^{-1}T^{n}u = 1$ soit $T^{n}u = u$.

\emph{Lemme 2} \quad Si un élément de $H^1(\mathbb{Z}_n, G_n)$ est représenté
par un $\alpha \in G_n$ tel que $\alpha\,T\alpha \cdots T^{n-1}\alpha = 1$,
alors son image dans $H^1(\mathbb{Z}_m, G_m)$ (si $n \mid m$) est aussi
représentée par $\alpha \in G_n \subset G_m$. \quad \add{Trivial}
\note{l'énoncé est marqué d'un trait vertical à gauche}

Posons
\[
  H^1(T, G) = \varinjlim H^1(\mathbb{Z}_n, G_n)
\]
\note{sous la limite, deux lignes : « \ill{} l'indice » et « filtrant
\uncertain{multiplicativement} »}

On voit donc que $H^1(T, G)$ est l'ens.\ des classes d'objets $\alpha \in G$
tels qu'il existe un $n$ avec $\alpha\,T\alpha \cdots T^{n-1}\alpha = 1$,
\struck{et} on écrivant $\beta \sim \alpha$ s'il existe $u \in G$ tel que
$\beta = u^{-1}\alpha\,Tu$.

\emph{Lemme 3} \quad Si \emph{Pour} \struck{\ill{}} l'application
$x \mapsto x^{-1}Tx$ de $G$ dans $G$ \emph{surjective}, on a
$H^1(T, G) = \lbrace e \rbrace$

\quad résulte des lemmes 1 et 2

\emph{Lemme 4} \quad Soit $e \to G' \to G \to G'' \to e$ une suite exacte de
$\mathbb{Z}$-groupes ($T$ opère dans $G$ et normalise $G'$). Supposons que
pour tout \ill{} l'application \struck{\ill{} \ill{} \ill{} \ill{} \ill{}}
$x \mapsto x^{-1}\alpha\,Tx$ \struck{de $G'$ dans $G'$ soit surjective}
\add{\uncertain{contienne} $G'$}. Alors
\[
  H^1(T, G) \to H^1(T, G'') \quad \text{est \emph{injectif}}
\]
\note{la ligne sous « l'application » est biffée et surchargée, avec
« intérieur » et « $G'$ » en interligne ; lecture de l'hypothèse
incertaine}

Soient en effet $\alpha, \beta \in G$, tels que $\alpha\,T\alpha \cdots
T^{n-1}\alpha = 1$ et de même $\beta\,T\beta \cdots T^{n-1}\beta = 1$, et

\page{79}

supposons qu'il existe $u \in G$ tel que $\beta \equiv u^{-1}\alpha\,Tu$
($G'$). On veut prouver que $\beta \sim \alpha$. Remplaçant $\alpha$ par
$u^{-1}\alpha\,Tu$, on peut supposer déjà que $\beta \equiv \alpha$ ($G'$)
i.e.\ $\beta = t\alpha$ ($t \in G'$). On cherche $v \in G'$ tel que
$\beta = v^{-1}\alpha\,Tv$ i.e.\ $t\alpha = v^{-1}\alpha\,Tv$ i.e.\
$t = v^{-1}\alpha\,Tv\,\alpha^{-1}$ ou encore, ($\sigma$ désignant l'aut.\
de $G'$ induit par l'aut.\ intérieur défini par $\alpha$)
$t = v^{-1}\sigma T(v)$. Un tel $v$ existe par hypothèse.

On cherche maintenant des conditions moyennant lesquelles
$H^1(T, G) \to H^1(T, G'')$ est surjectif. Soit donc $\alpha \in G$ tel que
$\alpha\,T\alpha \cdots T^{n-1}\alpha \equiv 1$ ($G'$), on veut montrer que
l'on peut trouver $t \in G'$ tel que $\beta = t\alpha$ satisfasse à
$\beta\,T\beta \cdots T^{n-1}\beta = 1$. Or le premier membre est
$(t\alpha\,T(t)T(\alpha) \cdots T^{n-1}(t)T^{n-1}(\alpha))$
\begin{align*}
  &= t\,\bigl(\alpha\,T(t)\,\alpha^{-1}\bigr)
  \bigl(\alpha\,T\alpha\,T^{2}(t)\,T(\alpha)^{-1}\alpha^{-1}\bigr) \cdots \\
  &\qquad \bigl(\alpha\,T\alpha \cdots T^{n-2}\alpha\,T^{n-1}(t)\,T^{n-2}\alpha^{-1}
  \cdots T\alpha^{-1}\alpha^{-1}\bigr)\,\alpha\,T\alpha \cdots T^{n-1}\alpha
\end{align*}
\note{chaque parenthèse est soulignée d'une accolade}
Le dernier facteur est $1$, d'où on trouve
\[
  t\,\sigma_1 T(t)\,\sigma_2 T^{2}(t) \cdots \sigma_{n-1}T^{n-1}(t)
\]
où $\sigma_i$ est l'automorphisme intérieur défini par
$\alpha\,T\alpha \cdots T^{i-1}\alpha$. D'où le \struck{\ill{}}

\emph{Lemme 5} \quad Soit $e \to G' \to G \to G'' \to e$ une suite exacte de
$\mathbb{Z}$-groupes, supposons que pour toute suite d'automorphismes
$\sigma_1, \ldots, \sigma_{n-1}$ de $G'$, induits par des automorphismes
intérieurs de $G$, l'application
\[
  x \longmapsto x\,\sigma_1 Tx \cdots \sigma_{n-1}T^{n-1}x
\]
de $G$ dans $G$ est surjective. Alors
\[
  H^1(X, G) \to H^1(X, G'')
\]
\note{la page s'arrête ici ; « $X$ » est tel qu'écrit, là où l'on attend
$T$}

\page{80}

Soit $\Gamma$ un groupe, $G$ un $\Gamma$-groupe, $M$ un $\Gamma$-\struck{groupe
\ill{}} \add{ensemble}. Supposons que $G$ opère sur $M$, et qu'on ait
\[
  \gamma(gm) = (\gamma g)(\gamma m)
\]
Soit alors $f$ un $1$-cocycle de $\Gamma$ à valeurs dans $G$
\[
  f(\gamma\gamma') = f(\gamma).\gamma f(\gamma')
\]
La donnée de $f$ permet de « tordre » $M$, de façon précise de définir une
autre façon d'opérer de $\Gamma$ sur $M$, savoir
\[
  \gamma_f\, m = f(\gamma).\gamma.m
\]
\note{il écrit $\gamma_f$ avec une petite flèche au-dessus du $\gamma$ et le
$f$ en dessous}
\begin{align*}
  [\text{car} \quad \gamma_f(\gamma'_f m) &= f(\gamma)\gamma\bigl(f(\gamma')
    \gamma' m\bigr) = f(\gamma)\,\gamma f(\gamma')\,\gamma\gamma' m \\
  &= f(\gamma\gamma')\,\gamma\gamma' m = (\gamma\gamma')_f\, m \\
  e_f\, m &= f(e).e.m = m \ ]
\end{align*}
On désigne par $M_f$ le nouveau \struck{\ill{}} \add{$\Gamma$-ensemble} ainsi
construit ; en fait, \uncertain{on ne s'est pas servi} de la loi de groupe de
$M$ ; en particulier, faisant opérer $G$ sur lui-même par automorphismes
intérieurs, on \struck{trouve sur $G_f$ une \ill{}} note que
\[
  \gamma(\sigma(g).g') = \sigma(\gamma g)\,\gamma g'
\]
car le \struck{\ill{}} \add{premier} membre est $\gamma(gg'g^{-1}) = \gamma g\,
\gamma g'\,(\gamma g)^{-1}$, c'est à dire le second membre. Ainsi, posant
\[
  \gamma_f\, g = f(\gamma)\,\gamma g\, f(\gamma)^{-1}
\]
on obtient sur $G$ une nouvelle loi de $\Gamma$-groupe.

Notons aussi \ill{} \ill{}
\[
  \gamma_f(gm) = \gamma_f(g)\,\gamma_f(m)
\]
car le \struck{deuxième} \add{premier} membre est
\[
  f(\gamma)\,\gamma(gm) = f(\gamma)\,\gamma g\,\gamma(m)
\]
et le deuxième
\[
  f(\gamma)\,\gamma g\, f(\gamma)^{-1} f(\gamma)\,\gamma(m)
\]
i.e.\ pareil.

\end{document}
